\chapter{Conclusions and Future Work}
\label{chap:conclusions}

This thesis investigated the operational feasibility, performance boundaries, and systemic trade-offs of centralized mathematical optimization for computational task offloading and multi-hop routing in Low Earth Orbit (LEO) satellite networks[cite: 4, 5, 9]. By designing, formalizing, and implementing a Centralized Integer Linear Programming (ILP) Orchestrator coupled with an event-driven dynamic batching mechanism within a \texttt{SimPy} discrete-event simulation environment, this work systematically explored the viability of orbital edge intelligence under realistic spaceborne constraints[cite: 3, 4].

This concluding chapter synthesizes the primary findings and answers the core research question formulated at the outset of the thesis[cite: 9]. Furthermore, it details the critical physical and architectural limitations inherent to centralized space computing and outlines promising future research directions toward resilient, hybrid orbital orchestration paradigms[cite: 5, 7].

% =============================================================================
% SECTION 6.1: SYNTHESIS OF CONTRIBUTIONS AND RESEARCH QUESTION RESOLUTION
% =============================================================================
\section{Synthesis of Contributions and Research Question Resolution}
\label{sec:concl_synthesis}

\subsection{Answering the Core Research Question}
\label{subsec:concl_rq_resolution}

In Section~\ref{sec:intro_objectives}, this thesis posed the foundational research question: \textit{To what extent can a centralized optimization approach maintain global allocation efficiency in high-density satellite networks before queuing and solver overhead cause systemic performance degradation?}[cite: 9]

Based on the multi-tiered experimental evaluations conducted across Chapters~\ref{chap:performance_analysis} and~\ref{chap:comparative_benchmark}, the answer is clear: **the performance ceiling of centralized orbital orchestration is fundamentally architectural rather than computational**[cite: 14].

\begin{itemize}
    \item \textbf{Computational Tractability of the Solver:} Contrary to the conventional assumption that NP-hard mathematical optimization is prohibitively slow for real-time spaceborne execution, the open-source \texttt{HiGHS} branch-and-cut solver demonstrated remarkable computational efficiency[cite: 3, 4]. By bounding the active optimization batch size ($BS \le 20$), the decision variable space remained strictly constrained, maintaining pure solver execution times at $T_{exec} \le 0.30\text{--}0.60\text{ ms}$ across all tested arrival rates ($\lambda \in [2, 10]\text{ req/s}$)[cite: 3, 4, 14]. This proves that solving regional linear formulations is computationally viable on contemporary space-qualified onboard processing units[cite: 3, 11].

    \item \textbf{Ingress Queuing as the Systemic Bottleneck:} System throughput collapse under elevated arrival rates ($\lambda \ge 6\text{ req/s}$) is governed almost entirely by the ingress queuing latency ($T_{wait}$) accumulated within the orchestrator buffer prior to solver invocation[cite: 3, 14]. As traffic scales, request aggregation delays severely erode task slack margins ($\Delta t_{rem} = D_r - t_{queue}$), triggering cascade \textit{Deadline Exceeded} dropouts before the scheduling algorithm can be executed[cite: 3, 14].

    \item \textbf{Primacy of Fast-Dispatch Policies:} To prevent queue-induced deadline collapse, centralized orchestrators must adopt responsive dispatch mechanisms[cite: 3]. Operating with an aggressive timeout ($BS = 20, BT \le 0.05\text{ s}$) maintains near-zero buffer residence times, effectively transforming the batching pipeline into a semi-continuous dispatch engine that sustains competitive completion rates ($74\%\text{--}78\%$ at low traffic, $40\%\text{--}44\%$ under stress)[cite: 3, 14].
\end{itemize}

\subsection{Algorithmic Archetypes and Workload Sensitivity}
\label{subsec:concl_workload_sensitivity}

The comparative benchmark in Chapter~\ref{chap:comparative_benchmark} established distinct operational envelopes across heterogeneous application profiles[cite: 3, 14]:

\begin{itemize}
    \item \textbf{Computation-Heavy Workloads ($\boldsymbol{\beta} = (0, 1, 0)$):} Distributed online heuristics (\texttt{OrbitAware}, \texttt{DTS-base}, \texttt{DTS-Optimal}) demonstrate superior resilience, achieving $72\%\text{--}74\%$ success rates at peak load ($\lambda = 10\text{ req/s}$)[cite: 5, 14]. Crucially, the rejection diagnostics revealed a complete operational divide[cite: 3]: heuristics dropped $100\%$ of rejected tasks due to physical capacity exhaustion (\textit{Insufficient CPU}), whereas ILP models suffered $60\%\text{--}88\%$ of their failures due to \textit{Deadline Exceeded}[cite: 3, 14]. In compute-bound environments, instantaneous local assignment without ingress buffering is essential to prevent temporal expiration[cite: 3, 14].

    \item \textbf{Communication-Heavy Workloads ($\boldsymbol{\beta} = (0, 0, 1)$):} The centralized optimization paradigm achieves clear superiority in data-intensive environments[cite: 14]. When offloading large datasets ($20\text{--}50\text{ MB}$), global visibility of multi-hop Inter-Satellite Link (ISL) states enabled the ILP formulations to attain near-perfect completion rates ($98\%\text{--}100\%$) under low-to-moderate loads, outperforming greedy heuristics[cite: 4, 14]. Furthermore, data-intensive workloads exhibited strong elasticity with respect to satellite battery capacity: expanding battery reserves from $40\text{ kJ}$ to $80\text{ kJ}$ yielded a $+15\%\text{--}25\%$ throughput recovery, proving that mathematically coordinated routing paths preserve energy across optical transceivers[cite: 3, 14, 17].

    \item \textbf{Multi-Objective Weight Calibration ($\alpha^* = 1000, \gamma^* = 8$):} The micro-tuning campaign proved that dominating the primary objective ($\alpha^* = 1000$) ensures that solver decisions are driven strictly by deadline compliance and capacity limits, yielding an average $+21\%$ gain in success rate[cite: 3]. Simultaneously, calibrating the orbital eclipse penalty to $\gamma^* = 8$ completely eliminated \textit{Sunset} drops ($0\%$ loss across all runs) while avoiding the $55\%$ latency inflation and topological path-bloat that occurs when the eclipse penalty is over-weighted ($\gamma \ge 1000$)[cite: 3, 6].
\end{itemize}

% =============================================================================
% SECTION 6.2: CRITICAL LIMITATIONS OF CENTRALIZED ORBITAL ORCHESTRATION
% =============================================================================
\section{Critical Limitations of the Centralized Approach}
\label{sec:concl_limitations}

Despite the mathematical rigor and global optimality demonstrated by the proposed framework, deploying a pure centralized orchestrator in operational mega-constellations involves severe structural limitations[cite: 5, 7].

\subsection{Ingress Buffering and the Batching Dilemma}
\label{subsec:limitations_batching}

Centralized mathematical scheduling requires aggregating incoming discrete requests into batches to instantiate the linear program[cite: 3, 4]. This creates an inescapable trade-off between optimization horizon and buffering latency[cite: 3, 9]:
\begin{equation}
    T_{total}(r) = T_{wait}(BS, BT, \lambda) + T_{up}(r, i) + W_{cpu,i} + \frac{d_{cpu,r}}{C_{sen,i}} + T_{down}(i, AP_r)
    \label{eq:concl_latency_chain}
\end{equation}
While larger batch sizes ($BS \ge 8$) allow the branch-and-cut algorithm to uncover superior combinatorial placement patterns across the fleet, the static wait time $T_{wait}$ directly consumes the application deadline slack[cite: 3, 4]. For mission-critical real-time services characterized by tight latency bounds ($D_r \le 1.0\text{ s}$), any artificial queuing introduced before solver execution severely degrades service availability[cite: 3, 9].

\subsection{Single Point of Failure and Space Hazards}
\label{subsec:limitations_spof}

Relying on a designated orbital Master node introduces an acute Single Point of Failure (SPOF)[cite: 5, 7]. Spaceborne edge computing operates in an exceptionally harsh environment subject to extreme thermal cycling, total ionizing dose (TID), and single-event upsets (SEUs) induced by cosmic rays and solar particle events[cite: 7, 9]. If the central Master satellite suffers a radiation-induced latch-up, an onboard computer reboot, an attitude control anomaly, or requires an autonomous orbital collision avoidance maneuver, the entire task scheduling and forwarding pipeline across the constellation is immediately paralyzed[cite: 5, 7].

\subsection{State Telemetry Overhead and Information Staleness}
\label{subsec:limitations_staleness}

To formulate an accurate mathematical representation of constellation constraints, the orchestrator must maintain an updated global state matrix containing processor queue lengths ($W_{cpu,i}$), residual battery levels ($B_i$), and active ISL bandwidths ($B_{u,v}$)[cite: 4]. Transmitting this telemetry across dynamic inter-satellite links incurs control signaling overhead that scales with the number of satellites ($\mathcal{O}(|S|)$)[cite: 4, 5]. 

Due to the extreme orbital velocities of LEO spacecraft ($\sim 7.5\text{ km/s}$), continuous relative motion alters line-of-sight angles and link propagation delays[cite: 4, 5]. Consequently, by the time state telemetry is gathered, buffered, and solved, the physical state of the constellation may have drifted[cite: 5]. This state staleness risks rendering calculated schedules sub-optimal or infeasible mid-execution, necessitating reliance on emergency greedy fallback routines[cite: 4].

% =============================================================================
% SECTION 6.3: FUTURE RESEARCH DIRECTIONS
% =============================================================================
\section{Future Research Directions}
\label{sec:concl_future_work}

The conclusions and limitations identified in this thesis pave the way for several natural extensions and advanced research trajectories[cite: 4, 5].

\subsection{Two-Tier Hybrid Orchestration Architectures}
\label{subsec:future_hybrid}

To overcome the trade-off between centralized optimality and distributed responsiveness, future architectures should investigate two-tier hierarchical frameworks[cite: 5]:
\begin{itemize}
    \item \textbf{Tier 1 (Fast-Path Heuristic Layer):} Time-critical and CPU-intensive requests are evaluated immediately at the local Access Point using decentralized, zero-queuing heuristics (such as \texttt{OrbitAware} or \texttt{DTS-Optimal}), ensuring sub-millisecond dispatching[cite: 5, 8].
    \item \textbf{Tier 2 (Global Optimization Layer):} Computationally relaxed, data-intensive tasks (such as delay-tolerant Earth Observation image filtering) are buffered and offloaded to localized cluster-head ILP solvers[cite: 4, 5]. This dual-track paradigm harnesses mathematical efficiency for bandwidth-heavy traffic while preserving ultra-low latency for time-sensitive services[cite: 3, 5].
\end{itemize}

\subsection{Traffic-Aware Adaptive Batching via Reinforcement Learning}
\label{subsec:future_adaptive_batching}

The static batching thresholds ($BS, BT$) evaluated in this work could be replaced by an intelligent, dynamic batching engine[cite: 3, 4]. A lightweight Reinforcement Learning (RL) agent—such as a contextual bandit or Q-learning controller—could be deployed at the ingress gateway to dynamically modulate the timeout threshold $BT(t)$ in real time[cite: 4, 5]:
\begin{equation}
    BT(t) = f\left( \hat{\lambda}(t), \, \min_{r \in Q} \{D_r - t_{arr,r}\}, \, \bar{W}_{cpu}(t) \right)
    \label{eq:adaptive_bt}
\end{equation}
By tracking traffic arrival surges ($\hat{\lambda}$) and the minimum residual slack time of buffered requests, the orchestrator could trigger immediate flushes during high-risk intervals and expand batch sizes during traffic lulls, optimizing the trade-off between scheduling efficiency and deadline preservation[cite: 3, 4].

\subsection{Multi-Agent Distributed Reinforcement Learning (MADRL)}
\label{subsec:future_madrl}

To eliminate the centralized Single Point of Failure and bypass the communication overhead of global telemetry aggregation, decentralized Multi-Agent Deep Reinforcement Learning (MADRL) presents a compelling alternative[cite: 5, 7]. Equipping each Satellite Edge Node with an autonomous actor-critic policy trained under Centralized Training with Decentralized Execution (CTDE) would allow satellites to make collaborative routing and scheduling decisions based solely on localized, one-hop neighborhood observations[cite: 4, 5]. This model could approximate the global efficiency of ILP solutions while achieving the instantaneous decision speeds characteristic of distributed heuristics[cite: 4, 5].

\subsection{Hardware-in-the-Loop (HIL) Prototyping on Space-Qualified Hardware}
\label{subsec:future_hil}

Finally, extending this simulation-based validation toward physical implementation requires testing the optimization pipeline on physical embedded hardware architectures[cite: 4]. Deploying the \texttt{HiGHS} solver and the batching engine on radiation-hardened system-on-chip (SoC) architectures—such as the ARM Cortex-R52, Cortex-A53, or fault-tolerant space-grade RISC-V platforms (e.g., NOEL-V)—will provide empirical measurements of CPU instruction overhead, memory footprint, cache contention, and power consumption during real-time solver execution[cite: 4, 11].

% =============================================================================
% SECTION 6.4: CONCLUDING REMARKS
% =============================================================================
\section{Concluding Remarks}
\label{sec:concl_closing}

The transition from passive bent-pipe relaying to active Orbital Edge Computing marks a significant evolution in space-based information systems[cite: 4, 5]. This thesis demonstrated that centralized Integer Linear Programming is computationally viable on spaceborne processing platforms, effectively bounding solver execution latency within sub-millisecond limits[cite: 4, 11, 14]. However, the temporal latency imposed by request buffering in ingress queues represents a fundamental physical boundary that limits the scalability of centralized scheduling under high-frequency traffic[cite: 3, 14]. 

By quantifying the operational trade-offs between centralized mathematical optimization, distributed hierarchical solvers, and autonomous heuristics across diverse mission profiles, this work provides architectural guidelines for designing the next generation of resilient, energy-aware orbital computing systems[cite: 3, 5, 14].